\section{Meet in the Middle}

Use Meet in the Middle (MITM) when $N$ is too large for $O(2^N)$ but $O(2^{N/2})$ is feasible---typically $N \approx 30$--$45$. Split items into two halves, enumerate both state sets, then combine them with sorting, binary search, or two pointers. For $N=40$, this reduces about $2^{40}$ states to two lists of $2^{20}$ states.

\begin{lstlisting}
using ll = long long;
void generate(const vector<ll>& a, int i, ll sum, vector<ll>& out) {
    if (i == (int)a.size()) { out.push_back(sum); return; }
    generate(a, i + 1, sum, out);
    generate(a, i + 1, sum + a[i], out);
}

// Exact subset sum target X:
vector<ll> L, R;
generate(leftHalf, 0, 0, L);
generate(rightHalf, 0, 0, R);
sort(R.begin(), R.end());
bool found = false;
for (ll x : L) found |= binary_search(R.begin(), R.end(), X - x);
\end{lstlisting}

For maximum subset sum $\leq X$, use \texttt{upper\_bound(R.begin(), R.end(), X-x)}, step back once, and maximize the pair sum. Sort both lists and use two pointers when suitable. To choose exactly $K$ elements or minimize a partition difference, store the selected count with each sum. MITM also applies to xor and other independently combinable half-states. Typical time is $O(2^{N/2} \cdot N)$ and memory is $O(2^{N/2})$.

\section{Pigeonhole principle}

If more than $N$ objects are assigned to $N$ boxes, two objects share a box. More generally, $N$ objects in $K$ categories guarantee one category with at least $\lceil N/K \rceil$ objects. The contest skill is choosing the boxes: remainders, prefix states, colours, or states in a finite process.

\subsection{Divisible subarray}
For $n$ integers, a non-empty contiguous subarray whose sum is divisible by $n$ always exists. Track prefix-sum remainders modulo $n$: a zero remainder gives a divisible prefix, and equal remainders give a divisible difference of prefixes.

\begin{lstlisting}
vector<int> first(n, -1);
first[0] = 0;
int rem = 0;
for (int i = 0; i < n; i++) {
    rem = (rem + a[i]) % n;
    if (rem < 0) rem += n;
    if (first[rem] != -1) {
        // Subarray [first[rem], i] has sum divisible by n.
        break;
    }
    first[rem] = i + 1;
}
\end{lstlisting}

Useful patterns: two values with difference divisible by $K$ (same remainder), prefix-state collisions, guaranteed duplicates, frequency lower bounds, and eventual cycles in a deterministic finite-state process. Ask: what are the objects, what are the possible states, and are there more objects than states?

\section{Kosaraju's algorithm}

Kosaraju finds strongly connected components (SCCs) in a directed graph. Vertices in one SCC mutually reach each other. Contracting every SCC produces the condensation graph, which is always a DAG---ideal for topological DP, longest paths, reachability, and source/sink analysis.

\begin{lstlisting}
struct Kosaraju {
    int n; vector<vector<int>> g, rg; vector<int> vis, order, comp;
    Kosaraju(int n) : n(n), g(n), rg(n), vis(n), comp(n, -1) {}
    void addEdge(int u, int v) { g[u].push_back(v); rg[v].push_back(u); }
    void dfs1(int u) {
        vis[u] = 1;
        for (int v : g[u]) if (!vis[v]) dfs1(v);
        order.push_back(u);
    }
    void dfs2(int u, int id) {
        comp[u] = id;
        for (int v : rg[u]) if (comp[v] == -1) dfs2(v, id);
    }
    int findSCCs() {
        for (int u = 0; u < n; u++) if (!vis[u]) dfs1(u);
        reverse(order.begin(), order.end());
        int id = 0;
        for (int u : order) if (comp[u] == -1) dfs2(u, id++);
        return id;
    }
};
\end{lstlisting}

The steps are: DFS on the original graph and save finishing order; reverse every edge; then DFS the transpose in decreasing finishing-time order. Each second-pass DFS is one SCC. Time and memory are $O(V+E)$. Common mistakes are forgetting the transpose, using the wrong order, and recursive stack overflow on very large graphs.

If more than one SCC exists, the standard minimum edges needed to strongly connect the condensation DAG is \texttt{max(number of source SCCs, number of sink SCCs)}. Kosaraju and Tarjan have the same $O(V+E)$ complexity; Kosaraju is often easier to remember, while Tarjan needs only one DFS.

\section{Recognition and practice}
\begin{itemize}
\item $N \approx 40$ with subset sum, partition, xor, closest sum, or choose-$K$: \textbf{MITM}.
\item More objects than remainders/states/categories; prefix sum plus modulo; guaranteed collision or cycle: \textbf{pigeonhole}.
\item Directed graph with mutual reachability, cycles that must become a DAG, or component compression: \textbf{Kosaraju/Tarjan}. 
\end{itemize}

Practice MITM with exact/closest/counting subset sums and two-pointer combinations; pigeonhole with modular prefix sums and repeated states; and SCCs with connectivity, source/sink components, minimum added edges, condensation-DAG DP, and weighted paths after compression.
